% One-page letter / statement template. Matches the CV letterhead.
\documentclass[11pt, letterpaper]{article}
\usepackage[T1]{fontenc}
\usepackage[utf8]{inputenc}
\usepackage{charter}
\usepackage[top=0.55in, bottom=0.45in, left=0.8in, right=0.8in]{geometry}
\usepackage[dvipsnames]{xcolor}
\definecolor{linkblue}{HTML}{1155CC}
\usepackage[colorlinks=true, urlcolor=linkblue, linkcolor=linkblue]{hyperref}
\usepackage{parskip}
\pagestyle{empty}
\setlength{\parindent}{0pt}
\linespread{1.0}

\newcommand{\letterhead}{%
  \begin{center}
    {\large\bfseries Rodrigo França Chaves}\\[2pt]
    {\small
      Chicago, IL \enspace$\cdot$\enspace
      \href{mailto:rchaves@uchicago.edu}{rchaves@uchicago.edu} \enspace$\cdot$\enspace
      +1 312 217 5526 \enspace$\cdot$\enspace
      \href{https://rodrigofch7.github.io}{rodrigofch7.github.io}
    }
  \end{center}
  \vspace{2pt}
  {\rule{\linewidth}{0.8pt}}
  \vspace{10pt}
}

\begin{document}
\letterhead
\begin{center}{\bfseries\large Statement of Interest: Professor Liz Moyer}\\[2pt]
{\small Mapping Energy History}\end{center}
\vspace{6pt}

Of the projects listed, this is the one I am most immediately equipped to start on, because I have spent the past year building exactly the pipeline it calls for. At the University of Chicago's Mansueto Institute I worked on converting century-old fire insurance registers, cloth-bound volumes that had never been machine-readable, into structured research data. I built the process end to end: page segmentation, optical character recognition, and structured extraction into clean records, then linkage of the results to three historical censuses by street, and classification of establishments into hundreds of industries using natural language processing over firm names. The output is a record of industrial location across twentieth-century Chicago that did not previously exist. I return to that project in September 2026 to write the method up as a paper documenting the pipeline for other researchers working with archival sources.

That is the same problem as building a database of the early electricity industry, tracing the spread of transmission lines, or identifying the earliest oil refineries. The difficulties are the ones I already know: inconsistent naming across sources and decades, addresses that must be resolved against street layouts that have since changed, entity matching without stable identifiers, and the constant judgement of when an uncertain match should be recorded as a match at all. I am comfortable with all three of the contribution routes listed in the description. I can do the geospatial mapping, having built gridded spatial panels and worked with shapefiles and rasters in QGIS and geopandas. I can do the digital-history data science, which is what the OCR and NLP pipeline above is. I am also willing to do straightforward archival research, which I have found is usually where you learn what the automated extraction is quietly getting wrong.

My substantive interest in this is genuine rather than convenient. My published work is on utility regulation, specifically whether transferring water and wastewater systems to private operators changed health outcomes across Brazilian municipalities between 1998 and 2021. Working on that made me interested in the question underneath this project, which is what actually promotes or constrains an energy transition. Ownership structure, regulatory design, and the spatial economics of network infrastructure recur in both, and the historical record is the only place to observe a completed transition rather than one in progress.

Two practical notes. I noticed the description mentions matching students with foreign-language fluency to appropriate projects. I am a native Portuguese speaker, fluent in English, conversational in Spanish, and have basic German, so Lusophone and Hispanophone sources would be available to me if any strand of the work reaches beyond the United States. Second, I take seriously that this project follows multi-author conventions and offers co-authorship. I would treat that as a commitment to see a contribution through to publication quality rather than to the end of a term.

I work in Python and R, with Git and reproducible pipelines, and I am used to documenting a method so that someone else can run it.

\end{document}
